\section{자동동형의 궤도와 기약인수}
\label{sec:normal-orbits-factorization}

\begin{learninggoal}
정규확장에서 켤레원소를 자동동형의 궤도로 해석한다. 기약다항식의 인수들이 자동동형 아래에서 서로 켤레임을 보이고 궤도-안정자 공식을 분리차수와 연결한다.
\end{learninggoal}

\subsection{켤레집합은 자동동형의 궤도이다}

\begin{theorem}
$L/K$가 정규 대수적 확장이고 $\alpha\in L$라 하자. 그러면
\[
  \Aut_K(L)\cdot\alpha
  =\operatorname{Conj}_K(\alpha).
\]
즉 $\alpha$의 모든 $K$-켤레는 $L$ 안에 들어 있고, $L$의 어떤 $K$-자동동형으로 $\alpha$에서 도달할 수 있다.
\end{theorem}

\begin{proof}
자동동형은 최소다항식을 보존하므로 왼쪽 궤도는 켤레집합에 포함된다.

반대로 $\beta$가 $\alpha$의 $K$-켤레라고 하자. 대수적 폐포 $\Omega\supseteq L$의 어떤 $K$-자동동형 $\tau$가 $\tau(\alpha)=\beta$를 만족한다. 정규성에 의해 $\tau(L)=L$이므로 제한
\[
  \tau|_L\in\Aut_K(L)
\]
이고 $\beta$는 왼쪽 궤도에 속한다.
\end{proof}

\begin{corollary}
$L/K$가 정규이고 $\alpha\in L$이면
\[
  \#\bigl(\Aut_K(L)\cdot\alpha\bigr)
  =[K(\alpha):K]_{\mathrm s}.
\]
$L/K$가 유한이고 $G=\Aut_K(L)$라 하면 궤도-안정자 정리에 의해
\[
  [K(\alpha):K]_{\mathrm s}
  =\frac{|G|}{|\Aut_{K(\alpha)}(L)|}.
\]
\end{corollary}

\begin{proof}
궤도는 최소다항식의 서로 다른 근들의 집합이므로 그 크기는 단순확장의 분리차수이다. 안정자는 $\alpha$와 $K$를 고정하는 자동동형들, 즉 $K(\alpha)$를 점별로 고정하는 자동동형들이다.
\end{proof}

\begin{example}
$L=\Q(\sqrt2,\sqrt3)$에서 $\alpha=\sqrt2+\sqrt3$의 궤도는
\[
  \sqrt2+\sqrt3,\quad
  \sqrt2-\sqrt3,\quad
  -\sqrt2+\sqrt3,\quad
  -\sqrt2-\sqrt3
\]
이다. 네 원소가 모두 서로 다르므로 안정자는 항등자동동형 하나뿐이다.
\end{example}

\subsection{기약인수들 사이의 대칭}

$\sigma\in\Aut_K(L)$와 다항식
\[
  h(X)=\sum_i a_iX^i\in L[X]
\]
에 대해
\[
  h^\sigma(X):=\sum_i\sigma(a_i)X^i
\]
라 쓰자.

\begin{theorem}[기약인수의 켤레성]
$L/K$가 정규 대수적 확장이고 $f\in K[X]$가 일계수 기약다항식이라 하자. $f$의 $L[X]$에서의 서로 다른 일계수 기약인수들을
\[
  f_1,\dots,f_r
\]
라 하자. 그러면 임의의 $i,j$에 대해 어떤 $\sigma\in\Aut_K(L)$가 존재하여
\[
  f_j=f_i^\sigma
\]
이다. 특히 모든 $f_i$의 차수는 같다.
\end{theorem}

\begin{proof}
대수적 폐포 $\Omega\supseteq L$에서 $f_i$의 한 근을 $\alpha$, $f_j$의 한 근을 $\beta$라 하자. 두 원소는 모두 $K$ 위 최소다항식 $f$의 근이므로 $K$-켤레이다. 따라서 어떤 $K$-자동동형 $\tau\in\Aut_K(\Omega)$가
\[
  \tau(\alpha)=\beta
\]
를 만족한다. $L/K$가 정규이므로 $\sigma=\tau|_L$는 $L$의 $K$-자동동형이다.

$f_i$는 $\alpha$의 $L$ 위 최소다항식이다. 계수와 근에 $\tau$를 적용하면 $f_i^\sigma$는 $\beta$를 근으로 갖는 일계수 기약다항식이다. 따라서 $\beta$의 $L$ 위 최소다항식 $f_j$와 같고
\[
  f_j=f_i^\sigma
\]
이다. 자동동형은 차수를 바꾸지 않으므로 모든 인수의 차수가 같다.
\end{proof}

\begin{remark}
$f$가 비분리이면 같은 기약인수가 거듭제곱의 중복도로 나타날 수 있다. 위 정리는 중복도를 제거한 서로 다른 기약인수들의 집합에 관한 명제이다. 이 인수들은 여전히 하나의 자동동형 궤도를 이룬다.
\end{remark}

\begin{example}
$L=\Q(\sqrt2)$ 위에서
\[
  X^4-2=(X^2-\sqrt2)(X^2+\sqrt2)
\]
이다. $\sqrt2\mapsto-\sqrt2$인 $\Q$-자동동형이 두 이차인수를 서로 바꾼다.
\end{example}

\begin{checkpoint}
\begin{enumerate}[label=(\alph*)]
  \item 정규확장에서 $\alpha$의 자동동형 궤도 크기가 최소다항식 차수보다 작을 수 있는 경우를 설명하라.
  \item $\Q(\sqrt2,\sqrt3)$에서 $\sqrt2$의 궤도와 안정자를 구하라.
  \item $X^4-2$의 두 $\Q(\sqrt2)$-기약인수가 같은 차수를 갖는 이유를 자동동형으로 설명하라.
\end{enumerate}
\end{checkpoint}
